\chapter[Discussão]{Discussão}
\label{cap:discussao}

Este capítulo interpreta os resultados obtidos, discute limitações do trabalho e
aponta ameaças à validade dos achados.

\section{Análise dos resultados}

O resultado central é a \textbf{ausência de vantagem \emph{consistente} do
pré-treinamento clínico}. Para a extração de relações no SemClinBr, no espaço de
três classes, o BioBERTpt não produz ganho estável sobre o BERTimbau, nem mesmo na
classe \texttt{negation\_of}, em que a especialização de domínio seria mais
plausível. Convém precisar o que os testes sustentam, porque o quadro é menos
confortável que um empate. Na métrica-alvo, o veredito de \emph{empate} vale para a
semente $42$ e não se repete na semente $43$, em que o \emph{bootstrap} pareado
sobre o F1 de \texttt{negation\_of} acusa diferença significativa, e em favor do
\emph{encoder} clínico. O que \emph{não} se mantém entre as duas sementes é a
\emph{direção nessa classe}, uma vez que o sinal da diferença se inverte. A
instabilidade é da métrica-alvo, e não do padrão global de erro. O teste de McNemar,
que incide sobre o
conjunto dos candidatos, rejeita a hipótese nula nas duas sementes e aponta nas duas
para o \emph{encoder} geral ($540$ contra $695$ discordâncias na semente $42$; $442$
contra $624$ na $43$), assim como a acurácia. O que se mantém, além dessa direção
agregada, é a \emph{magnitude relativa}. Em ambas as sementes, a variação de um
mesmo \emph{encoder} entre sementes supera a diferença medida entre \emph{encoders}
distintos (Seção~\ref{sec:robustez}). É nesse sentido, e apenas nesse, que não há
ganho atribuível ao domínio de pré-treino. O motivo não é uma derrota do clínico,
que chega a ficar à frente na negação sob a semente $43$, mas o fato de a diferença
entre os dois na classe investigada não sobreviver à troca de uma semente. Esse
achado contrasta com a expectativa difundida, bem como com os ganhos relatados
para o reconhecimento de entidades nomeadas clínicas em português \cite{schneider_2020_biobertpt}, de que
modelos de domínio sempre se transfiram melhor.

Duas explicações são plausíveis e compatíveis com os dados.

A primeira é \textbf{linguística}. A negação clínica em português é fortemente
lexicalizada (``não'', ``sem'', ``nega'', ``ausência de'', além de adjetivos que
embutem negação morfologicamente, como ``afebril'' e ``anictérico'', conforme o
Apêndice~\ref{ap:lexico}). Esse sinal é frequente e regular o bastante para que um
\emph{encoder} de domínio geral, com vocabulário amplo do português, o capture sem
desvantagem consistente em relação a um clínico. A caracterização do corpus reforça
essa leitura, pois a mediana da distância entre as entidades de uma relação
\texttt{negation\_of} é de um único caractere (Seção~\ref{sec:corpus}), ou seja,
o gatilho de negação está praticamente
colado ao achado negado. Detectá-lo depende pouco de conhecimento clínico
especializado.

A segunda é \textbf{estrutural}. A tarefa, tal como formulada, depende mais da
estrutura local do par, isto é, da direção, da proximidade e da delimitação das
entidades, do que de conhecimento lexical especializado de domínio, e essa
estrutura é fornecida
igualmente aos dois \emph{baselines} pela representação de entrada compartilhada
(Seção~\ref{sec:markers}). A semelhança entre as matrizes de confusão
(Seção~\ref{sec:confusao}) é a evidência mais direta disso. Os dois modelos não
apenas chegam a números próximos, mas concentram o erro no mesmo lugar, a
fronteira entre \texttt{associated\_with} e \texttt{no\_relation}.

Um corolário prático, e diretamente relevante para o desenvolvimento do
RECLin-PT, é a \textbf{eficiência}. O BERTimbau iguala o BioBERTpt em macro-F1
médio ($0{,}695$ contra $0{,}694$) com aproximadamente $60\%$ dos parâmetros
($108{,}9$ contra $177{,}9$ milhões). Quando o desempenho é equivalente, o modelo
menor é preferível por custo de inferência e de memória, consideração concreta para
implantação em ambientes clínicos com recursos limitados. A esse argumento soma-se a
estabilidade, dado que a amplitude entre sementes do BERTimbau na métrica-alvo ($0{,}008$) é
cerca de dez vezes menor que a do BioBERTpt ($0{,}078$), o que torna seu desempenho
mais previsível em uma nova execução. Vale registrar a contrapartida. É justamente
essa dispersão que permite ao BioBERTpt superar o geral na semente $43$, de modo que
a preferência pelo \emph{encoder} menor se justifica por previsibilidade e custo,
não por superioridade média na métrica-alvo.

\section{A classe \texttt{associated\_with}}

A classe \texttt{associated\_with} permanece a mais difícil para os dois modelos
(F1 $\approx 0{,}43$), com alta cobertura e baixa precisão. Duas causas se somam.

A primeira é \textbf{conceitual}. Distinguir associação genuína de mera proximidade
textual é intrinsecamente ambíguo em texto clínico telegráfico, em que entidades
vizinhas frequentemente não mantêm relação anotada, como ilustra o próprio
exemplo de geração de candidatos (Quadro~\ref{quad:exemplo_candidatos}), em que
cinco dos seis pares ordenados possíveis são \texttt{no\_relation}.

A segunda é \textbf{metodológica}. A janela \texttt{max\_gap}$=25$ descarta
$9{,}11\%$ das relações \texttt{associated\_with} do conjunto de teste
(Seção~\ref{sec:teto_recall}), enquanto não descarta \emph{nenhuma} relação
\texttt{negation\_of}. O F1 dessa classe é, portanto, medido sob um teto de
\emph{recall} substancialmente mais baixo que o das demais. Como o teto é idêntico
para os dois modelos, ele não afeta a comparação, objeto deste trabalho, mas
qualquer leitura do valor absoluto de F1 de \texttt{associated\_with} precisa levá-lo
em conta.

\section{Limitações}

A limitação mais séria recai sobre a própria caracterização de empate. Ela vale
para a semente $42$; a comparação pareada refeita na semente $43$ acusa diferença
significativa, e em favor do BioBERTpt clínico, invertendo o sinal observado na
semente de referência. A conclusão prática, a ausência de ganho \emph{confiável}
atribuível ao pré-treinamento clínico, sobrevive a essa inversão, porque se apoia na
instabilidade da diferença e não em sua direção; mas qualquer afirmação sobre
\emph{qual} \emph{encoder} é superior depende inteiramente da semente e não está
sustentada por estes experimentos. Além disso, \textbf{duas sementes por modelo não
bastam} para estimar um intervalo de confiança sobre essa variação, e o desvio
reportado na Tabela~\ref{tab:robustez} deve ser lido como dispersão observada, não
como estimativa de variância. Fechar a questão exige um número maior de repetições
com agregação formal.

A tarefa pressupõe \textbf{entidades já anotadas} (\emph{gold}). A avaliação não
considera o erro propagado de uma etapa de reconhecimento de entidades nomeadas
automática, de modo que os resultados refletem a extração de relações isoladamente.
Em um cenário fim a fim, os valores absolutos seriam necessariamente menores.

Por fim, o achado é específico desta \textbf{formulação de três classes} e desta
\textbf{escala de modelo} (\emph{base}). Não se pode excluir que o pré-treinamento
clínico volte a importar em espaços de relação mais ricos, em modelos maiores ou em
cenários com menos dados de ajuste.

\section{Ameaças à validade}

Quanto à \textbf{validade interna}, o congelamento das partições com semente fixa e
verificação por \emph{hashes} SHA-256 (Apêndice~\ref{ap:sha256}), somado ao
particionamento em nível de documento, mitiga o vazamento de dados, a principal
causa de resultados não reproduzíveis em PLN biomédico. A paridade total de
implementação entre os dois \emph{baselines} (Capítulo~\ref{cap:prototipo}) garante
que a diferença medida reflita o \emph{checkpoint}, e não artefatos de código.
Restam, contudo, decisões de projeto que influenciam os números absolutos, das
quais a mais consequente é a janela \texttt{max\_gap}$=25$ na geração de pares
candidatos. Esse limite foi quantificado por partição e por classe
(Seção~\ref{sec:teto_recall}) e, portanto, é uma limitação \emph{documentada e de
magnitude conhecida}, e não uma fonte de incerteza não medida.

Uma segunda ameaça diz respeito à \textbf{independência das instâncias}. O
\emph{bootstrap} pareado trata os $19\,210$ candidatos como independentes, o que
não é estritamente verdadeiro, pois pares de um mesmo documento compartilham texto e
vocabulário, de modo que o intervalo de confiança reportado é provavelmente mais
estreito que o real. Essa ameaça atenua-se, mas não desaparece, pelo fato de a
comparação ser \emph{pareada}, já que ambos os modelos são avaliados exatamente sobre as
mesmas instâncias, e a dependência afeta os dois igualmente. Ainda assim, a
conclusão de empate na semente $42$ (e a de diferença significativa na
semente $43$) deve ser lida com essa ressalva.

Quanto à \textbf{validade externa}, o SemClinBr, apesar de multi-institucional e
multiespecialidade, representa um conjunto específico de instituições e estilos de
redação; a generalização para outros contextos clínicos e para outros espaços de
relação deve ser feita com cautela. As inconsistências de anotação conhecidas do
corpus, ainda que tratadas e registradas pelo \emph{parser}, também impõem um limite
superior ao desempenho alcançável, limite que se aplica igualmente aos dois
modelos comparados.
